\documentclass[12pt,a4paper]{article}

% =====================================================
% PACOTES
% =====================================================
\usepackage[utf8]{inputenc}
\usepackage[T1]{fontenc}
\usepackage[brazil]{babel}
\shorthandoff{"}
\usepackage{geometry}
\usepackage{graphicx}
\usepackage{fancyhdr}
\usepackage{titlesec}
\usepackage[table]{xcolor}
\usepackage{hyperref}
\usepackage{setspace}
\usepackage{enumitem}
\usepackage{newtxtext,newtxmath}
\usepackage[most]{tcolorbox}
\usepackage{float}

\geometry{
    top=2.5cm,
    bottom=3cm,
    left=1.5cm,
    right=1.5cm
}

\onehalfspacing

% =====================================================
% CORES PERSONALIZADAS
% =====================================================
\definecolor{azulTitulo}{RGB}{30,60,90}
\definecolor{azulObjetivos}{RGB}{30,60,90}
\definecolor{cinzaCodigo}{RGB}{245,245,245}
\definecolor{vermelhoCodigo}{RGB}{180,50,50}
\definecolor{mostarda}{RGB}{204,153,0}

% =====================================================
% ESTILIZAÇÃO DE BOXES (IFSC STYLE)
% =====================================================
\newtcolorbox{boxAtividade}[1][]{
    breakable,
    enhanced,
    colback=azulTitulo!5,
    colframe=azulTitulo,
    boxrule=0pt,
    arc=6pt,
    outer arc=6pt,
    boxsep=0pt,
    left=14pt,
    right=14pt,
    top=16pt,
    bottom=14pt,
    title=#1,
    coltitle=white,
    colbacktitle=azulTitulo,
    fonttitle=\bfseries,
    attach boxed title to top left={
        yshift=-3mm,
        xshift=6mm
    },
    boxed title style={
        sharp corners,
        boxrule=0pt
    }
}

\newtcolorbox{boxCodigo}{
    colback=cinzaCodigo,
    colframe=vermelhoCodigo,
    boxrule=0.8pt,
    arc=2pt,
    left=6pt,
    right=6pt,
    top=6pt,
    bottom=6pt,
    borderline={0.5pt}{0pt}{vermelhoCodigo,dashed}
}

\newtcolorbox{boxDestaque}{
    colback=mostarda!10,
    colframe=mostarda,
    boxrule=1pt,
    arc=4pt,
    left=8pt,
    right=8pt,
    top=8pt,
    bottom=8pt
}

% =====================================================
% CABEÇALHO E RODAPÉ
% =====================================================
\pagestyle{fancy}
\fancyhf{}
\fancyhead[L]{
    \IfFileExists{../../uc-topicos-avancados-latex/figures/garopaba_horizontal_marca2015_PNG.png}{
        \includegraphics[height=1.4cm]{../../uc-topicos-avancados-latex/figures/garopaba_horizontal_marca2015_PNG.png}
    }{IFSC - Garopaba}
}
\fancyhead[R]{
    \textbf{Guia de Revisão - AV1}\\
    UC – Tópicos Avançados em Informática\\
    Técnico Integrado em Informática
}
\fancyfoot[L]{\small\textit{Revisão de Conteúdos e Comandos}}
\fancyfoot[R]{\small Página \thepage}
\renewcommand{\footrulewidth}{2.4pt}
\setlength{\headheight}{45pt}

% =====================================================
% ESTILO DE SEÇÃO
% =====================================================
\titleformat{\section}
{\normalfont\Large\bfseries\color{azulTitulo}}
{\thesection}
{1em}
{}
[\vspace{0.3cm}\hrule]

% =====================================================
% DOCUMENTO
% =====================================================
\begin{document}

\begin{center}
    {\LARGE \textbf{Material de Revisão para Avaliação Teórica 01}}\\
    \vspace{0.2cm}
    \textit{Preparação para Prova - Conteúdo: Roteiros 01 ao 07}
\end{center}

\vspace{0.5cm}

\section{Introdução: O Pipeline de Desenvolvimento}

Nesta UC, não apenas programamos; nós construímos um fluxo profissional:
\begin{itemize}
    \item \textbf{Versionamento (Git):} Segurança e colaboração.
    \item \textbf{Modelagem (OO):} Organização da lógica de negócio.
    \item \textbf{Serviço Web (Flask):} Disponibilização dos dados para o mundo.
    \item \textbf{Infraestrutura (Docker):} Portabilidade e isolamento.
\end{itemize}

\section{Revisão de Funções e Lógica em Python}

Antes da Orientação a Objetos, reforçamos o uso de funções para modularizar o código.

\begin{boxAtividade}[Conceitos de Funções]
\begin{itemize}
    \item \textbf{Definição:} Usamos \texttt{def nome\_funcao(parametros):}.
    \item \textbf{Retorno:} Toda função que processa um dado deve usar \texttt{return} para entregar o resultado.
    \item \textbf{Escopo:} Variáveis dentro da função são locais e não existem fora dela.
\end{itemize}
\end{boxAtividade}

\section{Orientação a Objetos (OO)}

O coração do desenvolvimento moderno em Python.

\begin{boxCodigo}
\begin{verbatim}
class Carro:
    def __init__(self, marca, modelo):
        self.marca = marca
        self.modelo = modelo
        self.ligado = False

    def ligar(self):
        self.ligado = True
\end{verbatim}
\end{boxCodigo}

\begin{itemize}
    \item \textbf{Classe:} A planta (Carro). \textbf{Objeto:} O carro na garagem.
    \item \textbf{Atributos:} Características (marca, modelo).
    \item \textbf{Métodos:} Comportamentos (ligar, acelerar).
    \item \textbf{self:} Referência ao próprio objeto que está sendo manipulado.
\end{itemize}

\section{Flask: Desenvolvimento Backend}

Transformando funções Python em rotas acessíveis via navegador.

\begin{itemize}
    \item \textbf{Ambiente Virtual (venv):} Crucial para isolar dependências (\texttt{pip install flask}).
    \item \textbf{JSON:} O formato padrão para troca de dados em APIs.
    \item \textbf{Métodos HTTP:}
        \begin{itemize}
            \item \textbf{GET:} Para buscar informações (visível na URL).
            \item \textbf{POST:} Para enviar dados de cadastro (corpo da requisição).
        \end{itemize}
\end{itemize}

\section{Docker: A Revolução dos Containers}

\begin{boxAtividade}[Comandos que salvam vidas]
\begin{itemize}
    \item \texttt{docker ps}: Lista o que está rodando agora.
    \item \texttt{docker logs -f <nome>}: Acompanha os erros do servidor em tempo real.
    \item \texttt{docker exec -it <nome> bash}: Entra "dentro" do container para debugar.
    \item \texttt{docker-compose down}: Desliga e remove tudo de forma limpa.
\end{itemize}
\end{boxAtividade}

\begin{boxDestaque}
\textbf{Lembre-se:} O Dockerfile é a \textbf{planta} da imagem. O Docker Compose é o \textbf{mestre de obras} que sobe vários serviços (ex: App + Banco) de uma vez só.
\end{boxDestaque}

\end{document}
